% Test-case specification TC-0021 (CTest test TC-0021-midstmt).
\documentclass[letterpaper]{article}

\usepackage[letterpaper,margin=1in]{geometry}
\emergencystretch=3em

% Compile from the repository root so this path resolves to
% the pinned WSP submodule.
\input{wsp/testing/test-case-library.tex}

\begin{document}

\def\TCID{TC-0021}
\def\TCTitle{%
  The Mid\$ statement}
\def\TCPurpose{%
  Verify \texttt{\_\allowbreak{}\_\allowbreak{}vbaMidStmt*}: in-place
  overwrite of a String, a fixed-length String buffer and a Variant, with
  the error cases.}
\def\TCPriority{\TCPriorityHigh}

\def\TCRequirementRef{%
  REQ-0003}
\def\TCDesignRef{%
  Not applicable.}
\def\TCFeatureRef{%
  \texttt{tests/\allowbreak{}vbrun\_\allowbreak{}midstmt\_\allowbreak{}tests.\allowbreak{}cpp},
  \texttt{src/\allowbreak{}helpers\_\allowbreak{}language.\allowbreak{}cpp}}
\def\TCTestTechnique{%
  Requirements-based testing; boundary start and length values.}

\def\TCPreconditions{%
  The library and the test program are built
  (\texttt{scripts/\allowbreak{}Invoke-Build.\allowbreak{}ps1}) and the test
  is registered with CTest.}
\def\TCEnvironment{%
  Windows with the toolchain of
  \texttt{scripts/\allowbreak{}Invoke-Build.\allowbreak{}ps1}; the local
  gate builds and runs x64 Debug, CI builds and runs the other
  architectures. The library is built from this repository in
  \texttt{out/\allowbreak{}build}.}
\def\TCAssumptions{%
  Not applicable.}

\def\TCInputData{%
  Strings, a fixed-length buffer and Variants (also by reference and null)
  with the start at zero and past the end, a zero and a negative length, an
  empty replacement and an empty or null target; the byte form.}
\def\TCInitialState{%
  No state from earlier runs is needed; the test makes what it needs.}

\def\TCProcedure{%
  \begin{enumerate}
    \item CTest runs \texttt{vbrun\_\allowbreak{}midstmt\_\allowbreak{}tests}.
    \item Overwrite part of each target with the helpers.
    \item Check the resulting text and the errors raised.
    \item CTest records the exit code and, on failure, the program's output
      (\texttt{--output-on-failure}).
  \end{enumerate}}
\def\TCExpectedResult{%
  All checks pass; the program prints \texttt{midstmt ok} and exits 0.}
\def\TCPostConditions{%
  No file or registry state remains that later tests depend on.}

\TCTemplate

\end{document}
